\section{Randbedingungen}
\label{sec:randbedingungen}

Randbedingungen (\emph{Constraints}) schränken den Lösungsraum ein.
Im Gegensatz zu Architekturentscheidungen, die das Team bewusst trifft
und in ADRs\footnotemark{} festhält, sind Randbedingungen
\emph{vorgegeben} -- etwa durch den Auftraggeber, die Organisation
oder das regulatorische Umfeld -- und können vom Architekturteam
nicht einseitig geändert werden.

\begin{hinweis}[Historischer Entscheidungsstand]
Einige technische Randbedingungen beschreiben den \emph{initialen}
Entscheidungsraum, bevor konkrete Architekturentscheidungen (ADRs)
getroffen wurden. Wo ADRs existieren, die eine offene Option
konkretisieren (z.\,B.\ Cloud-Provider, Frontend-Framework,
Message Broker), wird dies im jeweiligen Eintrag vermerkt.
Die verbindliche Entscheidung liegt stets im zugehörigen ADR
(Kap.~9).
\end{hinweis}

\footnotetext{%
  \textbf{ADR -- Architecture Decision Record.}
  Ein ADR ist ein kurzes, strukturiertes Dokument, das eine einzelne,
  architekturrelevante Entscheidung samt Kontext, betrachteten
  Alternativen und Begründung festhält. Das Format geht auf
  Michael Nygard zurück
  (\url{https://cognitect.com/blog/2011/11/15/documenting-architecture-decisions}).%
}

% ──────────────────────────────────────────────
\subsection{Technische Randbedingungen}
\label{sec:randbedingungen-technisch}

\begin{tabularx}{\textwidth}{L{3.2cm} X}
  \rowcolor{tableheader}
  \color{white}\textbf{Randbedingung} &
  \color{white}\textbf{Hintergrund / Erläuterung} \\

  Webtechnologie &
  Browserbasierter Zugriff ohne lokale Installation;
  Unterstützung mobiler Browser ist erforderlich. \\

  \rowcolor{tablegrey}
  Cloud-Deployment &
  Public Cloud; kein On-Premise-Betrieb.
  \textit{Historischer Stand zum Entscheidungszeitpunkt: Azure oder AWS präferiert.}
  Die finale Provider-Wahl ist in ADR-010 dokumentiert:
  Microsoft Azure, Region Switzerland~North (Zürich). \\

  REST / JSON &
  Alle Backend-Dienste stellen RESTful-HTTP-APIs bereit.
  GraphQL ist für eine spätere Ausbaustufe vorgesehen
  (vgl.\ Kap.~14.3). \\

  \rowcolor{tablegrey}
  Containerisierung &
  Docker-Container, orchestriert via Kubernetes
  (K8s~1.28+). \\

  HTTPS / TLS &
  Sämtliche Client-Server-Kommunikation erfolgt über HTTPS
  mit TLS~1.3. \\

  \rowcolor{tablegrey}
  Datenbank / Messaging &
  PostgreSQL~15 für transaktionale Daten;
  RabbitMQ als Message Broker (Azure Service Bus als evaluierte Alternative --
  zurückgestellt, vgl.\ ADR-002). \\

  Authentifizierung &
  OAuth~2.0 / OpenID Connect über einen externen Identity
  Provider (z.\,B.\ Keycloak). Kein proprietäres
  Auth-System. \\

  \rowcolor{tablegrey}
  Programmiersprachen &
  Backend: Java~21 (Spring Boot~3) oder Kotlin.\newline
  Frontend: TypeScript mit Angular (entschieden, vgl.\ ADR-005;
  Vue.js als gleichwertige Alternative evaluiert und zurückgestellt). \\
\end{tabularx}

% ──────────────────────────────────────────────
\subsection{Organisatorische Randbedingungen}
\label{sec:randbedingungen-organisatorisch}

\begin{tabularx}{\textwidth}{L{3.2cm} X}
  \rowcolor{tableheader}
  \color{white}\textbf{Randbedingung} &
  \color{white}\textbf{Hintergrund / Erläuterung} \\

  Teamgrösse &
  Cross-funktionales Team von 5--8~Personen
  (inkl.\ QA-Engineering). \\

  \rowcolor{tablegrey}
  Entwicklungsprozess &
  Scrum mit 2-Wochen-Sprints. Die Definition of Done
  schliesst Unit- und Integrationstests ein. \\

  Zeitrahmen &
  MVP innerhalb von 6~Monaten; vollständiger Rollout
  nach 12~Monaten. \\

  \rowcolor{tablegrey}
  Budget (Cloud-Betrieb) &
  Cloud-Betriebskosten $\leq$~CHF\,800\,/\,Monat.
  Ein Puffer von 10\,\% (CHF\,80) wird empfohlen
  (vgl.\ Kap.~11, R12). \\

  Open Source &
  Bevorzugte Nutzung von Komponenten unter Apache-2.0-
  oder MIT-Lizenz. \\

 \rowcolor{tablegrey}
  Source Control &
  Git (GitHub). Trunk-Based Development\footnotemark{} mit
  kurzlebigen Feature-Branches und verpflichtenden
  Pull-Request-Reviews. \\
\end{tabularx}

\footnotetext{%
  \textbf{Trunk-Based Development} bezeichnet ein
  Versionsverwaltungsmodell, bei dem es nur einen einzigen
  Hauptzweig (\texttt{main}) gibt, der jederzeit auslieferbar
  sein muss. Änderungen entstehen auf kurzlebigen Arbeitszweigen
  (\emph{Feature-Branches}, Lebensdauer 1--3~Tage) und werden
  über sogenannte \emph{Pull Requests} zurückgeführt: Mindestens
  ein weiteres Teammitglied prüft dabei den Code auf Korrektheit
  und Einhaltung der Projektstandards (\emph{Code Review}).
  Die Zusammenführung erfolgt erst nach erfolgreicher Prüfung
  und bestandener automatisierter Testpipeline.%
}


% ──────────────────────────────────────────────
\subsection{Rechtliche und regulatorische Randbedingungen}
\label{sec:randbedingungen-recht}

\begin{tabularx}{\textwidth}{L{3.2cm} X}
  \rowcolor{tableheader}
  \color{white}\textbf{Randbedingung} &
  \color{white}\textbf{Hintergrund} \\

  DSGVO / revDSG &
  Personenbezogene Daten unterliegen der EU-DSGVO sowie
  dem Schweizer revDSG. Datenspeicherung erfolgt
  ausschliesslich in der EU bzw.\ der Schweiz. \\

  \rowcolor{tablegrey}
  Datensparsamkeit &
  Es werden nur die für den Kernbetrieb notwendigen
  Personendaten erhoben (Grundsatz der
  Datenminimierung). \\

  Recht auf Löschung &
  Personendaten müssen auf Anfrage innerhalb von
  30~Tagen gelöscht oder anonymisiert werden
  (Art.~17 DSGVO; Art.~32 revDSG). \\

  \rowcolor{tablegrey}
  Auftragsverarbeitung &
  Sämtliche Cloud-Provider und externen Dienstleister
  sind über einen Auftragsverarbeitungsvertrag (AVV)
  nach Art.~28 DSGVO zu binden. \\

  Protokollierung &
  Sicherheitsrelevante Aktionen müssen auditierbar
  protokolliert werden. Aufbewahrungsfrist: mindestens
  3~Jahre. \\
\end{tabularx}
